\clearpage\label{answers}
{\sans\bfseries\LARGE\color{ink} Answers}\par
\small
Try every puzzle first. Where an answer shows a calculation, check the steps, not just the result.

\subhead{Chapter 1}
\textbf{Multix on the corner.} Staring cannot bend light. Light from behind the corner travels in straight lines and is blocked by the wall. \textbf{Three cards.} A dot shows when the holes line up. It disappears when one card moves. \textbf{Keep box.} Kept: the direction light goes. Set aside: its colour and brightness.

\textbf{1.1} (a) One shadow, if the torch is exactly in line. (b) Each pencil blocks the same line of light, so the shadows fall on top of each other.\\
\textbf{1.2} (a) S (b) N, a reflection is light bouncing, not blocked. (c) S (d) N.\\
\textbf{1.3} (a) A straight line from the torch through the hole to the wall. (b) Up. The dot moves 10~cm up, twice the card's move, because the wall is twice as far as the card.\\
\textbf{1.5} (a) $8\times10\div40=2$~cm. (b) Bigger: the nearer the object, the bigger the picture.\\
\textbf{1.4} Lines of light from the ball to your eye run through the wall and are blocked.

\subhead{Chapter 2}
\textbf{Multix on bigger shadows.} No. Moving the wall as well as the puppet changes two things, so the test cannot say which made the shadow bigger. \textbf{Keep box.} Kept: the lines of light and the heights. Set aside: the fuzzy edges of the shadow.

\textbf{2.1} (a) $10\times60\div15=40$~cm. (b) $10\times60\div60=10$~cm.\\
\textbf{2.2} (a) The distance from the puppet to the torch. (b) The height of the torch. (c) Keep the torch in place and change only the puppet's distance.\\
\textbf{2.3} (a) 12~cm. (b) 1~cm. (c) A 1~cm miss is about the size of the other misses. Measure again before changing your trust.\\
\textbf{2.4} (a) B. (b) S. (c) B.\\
\textbf{2.5} (a) $18\times60\div20=54$~cm. (b) $18\times60\div40=27$~cm.\\
\textbf{2.6} 20~cm and 40~cm. The second is twice as tall.\\
\textbf{2.7} $10\times80\div20=40$~cm. It grew, from 30~cm.\\
\textbf{2.8} About 20~cm. About 40~cm.\\
\textbf{2.9} Light travels in straight lines from the torch past the hand to the wall. A hand nearer the torch blocks a wider spread of lines, so the lines are steeper and land higher and wider on the wall.

\subhead{Chapter 3}
\textbf{Multix on noon.} No. The sun is highest at solar noon, which is not always 12:00 on the clock. \textbf{Keep box.} Kept: the direction and height of the sun, shown as a shadow. Set aside: its brightness and colour.

\textbf{3.1} (a) 12:00. (b) $45-26=19$~cm. (c) The 14:00 shadow, 35~cm against 34~cm.\\
\textbf{3.2} West. At 3:00 in the afternoon in Canada the shadow points east.\\
\textbf{3.3} (a) M (b) N (c) A.\\
\textbf{3.4} The sun is south. A high sun makes a short shadow.\\
\textbf{3.6} Yes. The shadows are almost equal in length (45 and 46~cm), so the halfway line points north.\\
\textbf{3.7} No. The lengths are very different, so the two tips are not at matching times either side of noon.\\
\textbf{3.8} December. A long noon shadow means the sun is low, which happens in winter.\\
\textbf{3.9} $76\div30\approx2.5$, so about 3 when rounded to the nearest whole number.\\
\textbf{3.5} $22\div20=1.1$, so 40~cm gives 44~cm and 10~cm gives 11~cm.

\subhead{Chapter 4}
\textbf{Multix on the dark.} No. A clock only works if the dots mark where the sun's shadow really falls at each hour. \textbf{Keep box.} Kept: the shadow's place at marked hours. Set aside: minutes between marks and the season.

\textbf{4.1} 11:30. The friend is wrong: a long shadow could be morning or afternoon. The place of the tip is what tells the time.\\
\textbf{4.2} (a) 6, 9 and 3~min. (b) $18\div3=6$~min. (c) $9-3=6$~min.\\
\textbf{4.3} (a) $12-4=8$~min. (b) Yes, 8 is bigger than 2 and 3. (c) The clock drifts as the weeks pass.\\
\textbf{4.4} Question: does the shadow clock still agree with a watch? Change: the day. Measure: minutes between shadow clock and watch. Keep the same: the stick, the board, the place, the hours checked.\\
\textbf{4.5} Clouds block the sun, so there is no shadow. People used other tools, such as a water clock or a marked candle.

\subhead{Chapter 5}
\textbf{Multix on angles.} No. Measure again and check the protractor. A 30 degree beam leaves at 30, not 60. \textbf{Keep box.} Kept: the angle in and angle out. Set aside: the width of the beam and the surface's bumps.

\textbf{5.1} Maze 1: F (3 mirror bounces). Maze 2: N (5 bounces). Maze 3: K (6 bounces).\\
\textbf{5.2} North hitting (/) goes east. West hitting (\textbackslash) goes north.\\
\textbf{5.3} Two parallel mirrors, each at 45 degrees, one at the top of the periscope and one at the bottom.\\
\textbf{5.5} MOM, WOW and TOOT.\\
\textbf{5.6} 9:00. The mirror swaps left and right, so the hour hand moves to the other side. Reasons for the second part vary; any careful reason that mentions the mirror flipping the picture left to right is fine.\\
\textbf{5.7} Any maze with three mirrors, checked by tracing.\\
\textbf{5.8} Maze A leaves at R, with 5 bounces. Maze B leaves at D, with 5 bounces.\\
\textbf{5.9} Yes. Light paths can be run in reverse. Sent in at F, the beam of Maze 1 comes out where the first beam entered.\\
\textbf{5.4} 40~degrees. Off by 4 degrees. That is a bit more than the 1 to 2 degrees in the practice runs, so measure again.

\subhead{Chapter 6}
\textbf{Multix on the model.} No. A model keeps some things and drops others. A shadow clock drops the minutes, the season and the weather.

\textbf{6.1} (a) S (b) B (c) G.\\
\textbf{6.2} (a) M (b) R (c) M.\\
\textbf{6.4} 8:00 L, 10:00 M, 12:00 S, 14:00 M, 16:00 L.\\
\textbf{6.5} (a) The slope of the ground and the sticks' place. (b) Two things changed, so a bigger shadow cannot be traced to one.\\
\textbf{6.6} The building's shadow is 10 times the stick's shadow, so the building is 10 times as tall: 300~cm.\\
\textbf{6.7} Keeps the direction of the shadow, drops the minutes or the season.\\
\textbf{6.3} Examples: the colour of the sun, the minutes, the season, the weather.\\
\textbf{Slop check.} (a) You cannot tell: no title, no axes, no labels. (b) Many colours, stars, joke text, tilted filler. (c) A title, the hours, the shadow lengths and units. (d) The stars and jokes.\\
\textbf{Redraw.} One title, labelled axes with hours and centimetres, calm colours.

\textbf{Whole-book table.} Chapter 2: simplified a shadow to a line drawing, kept the straight lines and the heights, set aside fuzziness. Chapter 3: simplified the sun to a direction and a height, kept the shadow length and direction, set aside brightness. Chapter 4: simplified the day to marked dots, kept the hour dots, set aside the minutes and the season. Chapter 5: simplified a mirror to a rule, kept the angles, set aside bumps. Chapter 6: simplified a project to a report, kept the pattern, set aside the detail of each step.

\clearpage\label{notes}
{\sans\bfseries\LARGE\color{ink} Words}\par
\small
\begin{multicols}{2}
\word{gnomon}{the stick or edge on a sundial whose shadow tells the time.}
\word{model}{a simple stand-in that keeps some things and drops others.}
\word{normal}{a line at a right angle to a mirror at the point where light hits.}
\word{parallel}{lines that run side by side and never meet.}
\word{ratio}{how many times one number holds another, such as shadow divided by height.}
\word{reflect}{to bounce off a surface.}
\word{shadow}{the place where light is blocked.}
\word{solar noon}{the time when the sun is highest, and the shadow shortest.}
\word{sundial}{a clock that uses a shadow.}
\end{multicols}

\vspace{0.4em}
{\sans\bfseries\large\color{sky!70!black} Notes for grown-ups}\par
This book teaches that a shadow, a drawing and a clock are models: simple stand-ins for light and the sun. Each chapter has one activity, one worked example with made-up practice numbers, one trust box, one ``what did we keep?'' prompt and a page of puzzles.

\textbf{Materials and safety.} A torch, a stick, card, a mirror and a sunny day. Never look at the sun, and never aim a mirror at eyes or use sunlight on one. Keep the torch low. The shadow tracing is done on a flat, safe spot outdoors.

\textbf{Directions.} The shadow directions in Chapter~3 are for the northern half of the world. South of the equator, the midday shadow points south.

\textbf{The practice numbers} were made up so the arithmetic is easy to check. Real shadows depend on the place, the day of the year and the clock, so real results will differ.

\subhead{Sources}
\begin{itemize}[leftmargin=1.3em,itemsep=0.2em]
\item E. Hecht. \emph{Optics}, 5th edition. Pearson, 2016. For the law of reflection and straight-line light.
\item A. E. Waugh. \emph{Sundials: Their Theory and Construction}. Dover, 1973. For the gnomon and how sundials work.
\item Cleomedes. \emph{On the Circular Motions of the Celestial Bodies}. For the later report of Eratosthenes' shadow measurement.
\end{itemize}

\clearpage
\thispagestyle{blank}
\begin{tikzpicture}[remember picture,overlay]
  \fill[sun!35] ($(current page.south west)+(0.4in,0.4in)$) rectangle ($(current page.north east)-(0.4in,0.4in)$);
  \node[anchor=north,align=center,font=\sffamily\bfseries,text=ink,text width=5.4in] at ($(current page.north)+(0,-1.4in)$) {\Large Next in Science\\[6pt]\fontsize{28}{30}\selectfont Things That Push Back};
  \node[align=center,font=\sffamily,text width=4.6in,text=ink] at (current page.center) {A box slides across carpet and stops sooner than it does on paper. Book 4 shows how to find the pushes and pulls on a thing, draw each as an arrow, and balance a seesaw from a ruler.};
  \node at ($(current page.south)+(0,2.6in)$) {\scalebox{3}{\multixbot[ok]}};
  \node[font=\sffamily\bfseries,text=ink] at ($(current page.south)+(0,0.9in)$) {Ages 9 to 11 \quad$\star$\quad Science Book 3};
\end{tikzpicture}
